\subsection{Matched progress-score ablation at 75k interactions}
\label{sec:routing75}

To isolate the contribution of using the progress scores, we added
a score-neutral condition for the 20 original BAMOR seeds on each
main task. It computes the same statistics but routes with all scores
set to zero, retaining the same backbone, stochastic collection
rule, operator library, warmup and family quotas. Thus the treatment
is access to scores for allocation, not the cost of computing scores.
The score-on arm reuses the original runs. This is a retrospective
completion on the inspected cohort, with shared training mode histories,
evaluation weights and reset seeds.

Turning scores on changed mean direct utility by $-14.41$ on HalfCheetah
and $+15.95$ on Hopper. Both interval methods include zero for both tasks
(Table~\ref{tab:routing75}). The matched comparison does not establish
that using the scores improves final-policy utility.

% Generated by make_evidence_tables.py from completed readouts.
\begin{table}[!htbp]
\centering\footnotesize
\caption{Matched score-on versus score-neutral BAMOR at 75k, 20 paired seeds per task. Both arms retain the quota mechanism, stochastic collection and the same backbone. Retrospective completion with the original score-on records. Intervals are marginal 97.5\%; W/L/T counts wins, losses and ties for score-on.}
\label{tab:routing75}
\begin{tabular}{@{}lrrrrr@{}}
\toprule
Task & $U$: on / neutral & $\Delta U$ & Bootstrap & Paired $t$ & W/L/T \\
\midrule
HalfCheetah & $296.3/310.7$ & $-14.41$ & ${[-78.40,53.62]}$ & ${[-86.69,57.88]}$ & 6/9/5 \\
Hopper & $436.6/420.6$ & $+15.95$ & ${[-63.56,97.75]}$ & ${[-73.35,105.25]}$ & 10/9/1 \\
\bottomrule
\end{tabular}
\end{table}


Of the 46 complete blocks, the first ten are scheduled. The 36
subsequent complete blocks use the selector, whether it is correcting a
quota deficit or sampling a family with no deficit. The final 1,000-step
partial block is selected but emits no complete-block feedback.
Tables~\ref{tab:slot_shift} and~\ref{tab:slot_counts} report how block
counts and actual assignments change across the paired runs. On average,
the complete post-warmup slot differs at 12.9 of 36 positions on
HalfCheetah and 17.3 on Hopper. The five exact utility ties on HalfCheetah
and the one on Hopper are also the pairs with identical logged
assignments, including the partial block. Subsequent trajectory and
random-stream differences can contribute to the assignment differences,
so these are not one-step effects of a score.

\newcommand{\NcaTrainingInteractions}{3,000,000}
\newcommand{\NcaEvaluationInteractions}{8,209,456}
\newcommand{\NcaEnvironmentInitializationProbeSteps}{80}

The additional condition consumed \NcaTrainingInteractions{} training
interactions, \NcaEvaluationInteractions{} final-policy evaluation
interactions and \NcaEnvironmentInitializationProbeSteps{} initialization
probe steps.
\FloatBarrier
